\begin{proposition}[4.5.9]
\label{IV.2.4.5.9-fr}
Sea $X$ un $k$-preesquema. Las condiciones siguientes son equivalentes:
\begin{enumerate}
\item[\textup{a)}] $X$ es geométricamente irreducible \textup{(es decir, para \emph{toda}
extensión $K$ de $k$, $X\otimes_k K$ es irreducible)}.
\item[\textup{b)}] Para toda extensión finita separable $K$ de $k$, $X\otimes_k K$
es irreducible.
\item[\textup{c)}] $X$ es irreducible y, si $x$ es su punto genérico, $\boldsymbol{k}(x)$ es una
extensión primaria de $k$.
\end{enumerate}
\end{proposition}

\begin{proof}
Es claro que $a)$ implica $b)$. Para ver que $b)$ implica $c)$, consideremos
una extensión finita separable $K$ de $k$; si $p:X\otimes_k K\to X$ es la
proyección canónica, los

\oldpage[IV-2]{63}

puntos maximales de $X\otimes_k K$ son los de la fibra
$p^{-1}(x)=\operatorname{Spec}(\boldsymbol{k}(x)\otimes_k K)$ ((2.3.4) y
$(\mathbf{0}_{\mathrm I},2.1.8)$); decir que esta fibra es irreducible para toda
extensión finita separable $K$ de $k$ equivale a decir que $\boldsymbol{k}(y)$ es una
extensión primaria de $k$, en virtud de (4.3.2). Recíprocamente, si se cumple
$c)$, el mismo razonamiento (teniendo en cuenta (4.3.2)) muestra que, para \emph{toda}
extensión $K$ de $k$, $X\otimes_k K$ es irreducible; por tanto, $c)$ implica
$a)$.
\end{proof}

\begin{corollary}[4.5.10]
\label{IV.2.4.5.10-fr}
Sean $X$ un $k$-preesquema irreducible, $x$ su punto genérico, $k'$ la
clausura algebraica separable de $k$ en $\boldsymbol{k}(x)$ y $k''$ una extensión
galoisiana de $k$ (de grado finito o no) que contiene a $k'$. Se supone que
$k'$ es finito sobre $k$; entonces las componentes irreducibles de
$X\otimes_k k''$ son geométricamente irreducibles; su número es igual a
$[k':k]$ y es también el número geométrico de componentes irreducibles de
$X$.
\end{corollary}

\begin{proof}
Al igual que en (4.5.9), la cuestión se reduce a considerar los puntos maximales de la
fibra $\operatorname{Spec}(\boldsymbol{k}(x)\otimes_k k'')$; en virtud de (4.3.4), son
$[k':k]$ y los cuerpos residuales en estos puntos (que son también los puntos
maximales de $X\otimes_k k''$) son extensiones primarias de $k''$, lo que,
teniendo en cuenta (4.5.9), demuestra el corolario.
\end{proof}

\begin{corollary}[4.5.11]
\EGAReaderTarget{ega4.c04.cor.4-5-11}{IV.2.4.5.11-fr}
Sea $X$ un $k$-preesquema; supongamos que $X$ solo tiene un número finito de
puntos maximales $x_i$ $(1\leqslant i\leqslant r)$ y que, para cada $i$, la
clausura algebraica separable $k'_i$ de $k$ en $\boldsymbol{k}(x_i)$ es de grado finito
sobre $k$. Entonces existe una extensión finita separable $L$ de $k$ tal que
las componentes irreducibles de $X\otimes_k L$ son geométricamente
irreducibles; su número, igual a $\sum_i[k'_i:k]$, es el número geométrico de
componentes irreducibles de $X$.
\end{corollary}

\begin{proof}
Los puntos maximales de $X\otimes_k L$ son los de las distintas fibras
$\operatorname{Spec}(\boldsymbol{k}(x_i)\otimes_k L)$, en virtud de (2.3.4) y
$(\mathbf{0}_{\mathrm I},2.1.8)$ aplicado a las componentes irreducibles de $X$ y a
sus imágenes inversas en $X\otimes_k L$. Basta, pues, tomar como $L$ una
extensión galoisiana finita de $k$ que contenga todos los $k'_i$ y aplicar
(4.5.10).
\end{proof}

Se observará que (4.5.11) se aplica, en particular, a todo $k$-preesquema
$X$ \emph{de tipo finito sobre $k$}, pues $X$ es entonces noetheriano y, por tanto,
solo tiene un número finito de componentes irreducibles, y cada $\boldsymbol{k}(x_i)$ es
una extensión de tipo finito de $k$, por lo que también lo es la clausura
algebraica de $k$ en $\boldsymbol{k}(x_i)$ ([1], p.~6-06, lema~5).

\begin{remarks}[4.5.12]
\label{IV.2.4.5.12-fr}
\begin{enumerate}
\item[(i)] Las nociones definidas en (4.5.2) dependen del cuerpo base $k$.
Para abreviar, y cuando sea necesario mencionar el cuerpo base, podrá decirse
«~$k$-irreducible~» (resp. «~$k$-conexo~») en lugar de
«~geométricamente irreducible (resp. conexo) relativamente a $k$~». Se
observará que esta terminología, completamente natural en el contexto de los
esquemas, es exactamente opuesta a la de Weil: en este autor,
«~$k$-irreducible~» (resp. «~$k$-conexo~», «~$k$-normal~», etc.) designa una
propiedad \emph{intrínseca} de un preesquema $X$, independiente del cuerpo
$k$ tomado como «~cuerpo base~» (es decir, del morfismo
$X\to\operatorname{Spec}(k)$ elegido); expresamos estas propiedades diciendo
simplemente que $X$ es irreducible (resp. conexo, normal, etc.). Por otra
parte, Weil dice «~\emph{absolutamente irreducible}~», «~\emph{absolutamente conexo}~»,
«~\emph{absolutamente normal}~», etc., para las nociones correspondientes relativas
al preesquema $X\otimes_k\Omega$, donde $\Omega$ es una extensión
algebraicamente cerrada adecuada de $k$. Esta terminología se explica por el
distinto punto de vista adoptado por este autor: para él, una «~variedad
algebraica~» $V$ se da \emph{ante todo} como un objeto

\oldpage[IV-2]{64}

geométrico sobre un cuerpo \emph{algebraicamente cerrado} $\Omega$; el dato de un
«~cuerpo de definición~» $k$ más pequeño (es decir, de un esquema sobre $k$
que define $V$ por extensión del cuerpo base a $\Omega$) es para él una
estructura adicional y, en cierta medida, secundaria, de modo que su
calificación de «~absoluto~» o «~intrínseco~» por un lado, «~relativo a
$k$~» por otro, es opuesta a la nuestra\footnote{El punto de vista de Zariski
es ya más próximo al nuestro, y su terminología no suele entrar en conflicto
con la introducida aquí.}. Por tanto, evitaremos en lo sucesivo utilizar el
calificativo «~absoluto~», que puede prestarse a confusión, y cuando
utilicemos la terminología abreviada introducida anteriormente (en conflicto
con la terminología establecida), remitiremos a (4.5.12) para evitar toda
ambigüedad.

\item[(ii)] La proposición (4.5.9) da un criterio «~birracional~» (es decir,
que solo depende del cuerpo residual en el punto genérico) para que un
$k$-preesquema $X$ sea geométricamente irreducible. No existe un criterio
análogo para que $X$ sea geométricamente conexo. Por ejemplo, tomemos
$X=\operatorname{Spec}(\mathbf R[[T,U]]/(T^2+U^2))$ ($T,U$ indeterminadas);
$X$ es un $\mathbf R$-esquema íntegro y su cuerpo de funciones racionales
$\boldsymbol{k}(x)$ es $\mathbf C((T))$, como se comprueba fácilmente;
$X\otimes_{\mathbf
R}\mathbf C$ se descompone en dos componentes irreducibles
(las dos «~rectas isótropas~») que tienen un punto común (el ideal maximal,
imagen en $\mathbf C[[T,U]]/(T^2+U^2)$ del ideal maximal $(T)+(U)$ de
$\mathbf C[[T,U]]$); por tanto, $X\otimes_{\mathbf R}\mathbf C$ es conexo.
Sea $X'$ el complementario (abierto) del punto cerrado de $X$ correspondiente
al ideal maximal $(T)+(U)$ de $\mathbf R[[T,U]]$; entonces
$X'\otimes_{\mathbf R}\mathbf C$ no es conexo, aunque $X$ y $X'$ tengan el
mismo cuerpo de funciones racionales.
\end{enumerate}
\end{remarks}
